\subsection{Voting Rights Act and Section 2}

Section 2 of the VRA prohibits voting practices that ``result in a denial or abridgement of the right to vote on account of race.'' The 1982 amendments established that plaintiffs need not prove discriminatory intent; demonstrating discriminatory effect suffices~\citep{vra1982}.

\textbf{Thornburg v. Gingles} (1986) established three preconditions for Section 2 liability:
\begin{enumerate}
    \item The minority group is ``sufficiently large and geographically compact''
    \item The minority group is politically cohesive
    \item The majority votes sufficiently as a bloc to defeat minority-preferred candidates
\end{enumerate}

Our work addresses the first precondition: \emph{What constitutes ``sufficiently large''?}

\subsection{Gingles Framework and Study Scope}

The three Gingles preconditions establish a structured framework for evaluating Section 2 vote dilution claims. Understanding what each prong requires---and critically, which aspects can be assessed algorithmically---is essential for properly scoping our empirical analysis.

\subsubsection{The Three Gingles Preconditions}

\textbf{Prong 1: Sufficiently Large and Geographically Compact}

The minority group must be sufficiently numerous and geographically concentrated to constitute a majority in a reasonably compact single-member district. This prong has two components:

\begin{itemize}
    \item \textbf{Numerosity}: Enough minority population to form a majority ($>$50\%) in at least one district
    \item \textbf{Geographic compactness}: Minority population concentrated enough that a district containing them would satisfy traditional compactness requirements (contiguity, reasonable shape)
\end{itemize}

Courts assess geographic compactness using both quantitative metrics (Polsby-Popper scores, Reock scores, edge-cut minimization) and qualitative factors (respecting political subdivisions, maintaining communities of interest, following natural boundaries).

\textbf{Prong 2: Politically Cohesive}

The minority group must demonstrate political cohesion---voting as a bloc for its preferred candidates with sufficient consistency to constitute a politically unified community. This requires:

\begin{itemize}
    \item Analysis of past election results showing minority voters' preferred candidates
    \item Statistical demonstration that minority voters support these candidates at rates significantly higher than chance
    \item Evidence of cohesion across multiple elections and election types (presidential, congressional, state/local)
\end{itemize}

Courts typically use ecological regression, homogeneous precinct analysis, or ecological inference to estimate minority voting patterns from aggregate precinct-level data.

\textbf{Prong 3: Majority Bloc Voting}

White voters (or the majority group) must vote as a bloc sufficiently to usually defeat the minority's preferred candidates in the absence of special circumstances (like incumbency or unusual candidate characteristics). This requires:

\begin{itemize}
    \item Demonstration that white voters consistently support candidates opposed to minority-preferred candidates
    \item Evidence of racially polarized voting patterns---minority and majority groups supporting different candidates
    \item Analysis across multiple elections showing the pattern is persistent, not anomalous
\end{itemize}

This prong ensures that minority vote dilution actually occurs---that district configurations matter because majority voting patterns prevent minority electoral success.

\subsubsection{What This Study Addresses: Prong 1 Only}

Our algorithmic analysis addresses \textbf{only the geographic compactness component of Gingles prong 1}. Specifically, we examine:

\begin{itemize}
    \item \textbf{Numerosity question}: Given a state's minority percentage, can optimization create districts with $\geq$50\% minority population?
    \item \textbf{Proportionality question}: Can optimization create MM districts in proportion to state minority percentage?
    \item \textbf{Geographic feasibility}: Do demographic-geographic constraints prevent proportional MM representation even with optimal algorithms?
\end{itemize}

We operationalize "geographically compact" as districts that:
\begin{enumerate}
    \item Are contiguous (all parts connected)
    \item Minimize edge-cut (number of census tract boundaries crossed between districts)
    \item Maintain population equality within $\pm$0.5\% of target
\end{enumerate}

Edge-cut minimization serves as proxy for traditional compactness. Lower edge-cuts correlate with districts that respect political subdivisions, maintain communities of interest, and produce reasonably regular shapes.

\subsubsection{What This Study Does Not Address: Prongs 2-3}

Our analysis \textbf{cannot and does not address political cohesion or racially polarized voting}. These preconditions require:

\textbf{Election data}: Precinct-level voting results showing minority and majority candidate preferences. Our demographic analysis uses only Census population data.

\textbf{Statistical inference}: Ecological regression or homogeneous precinct analysis to estimate group voting patterns. These methods require electoral data and sophisticated statistical modeling beyond algorithmic redistricting.

\textbf{Case-specific context}: Political cohesion and racially polarized voting vary by jurisdiction, election type, candidate characteristics, and temporal period. They cannot be determined from demographic data alone.

\textbf{Behavioral analysis}: Prongs 2-3 concern voter behavior and electoral outcomes, not geographic configurations. Algorithmic redistricting addresses district shapes and demographic composition, not voting patterns.

\subsubsection{Implications for VRA Analysis}

This scope limitation has critical implications:

\textbf{Necessary but not sufficient}: Our findings establish whether geographic feasibility for proportional MM representation exists (Gingles prong 1). This is necessary for Section 2 liability but not sufficient. Even if proportional MM districts are geographically feasible, plaintiffs must still demonstrate political cohesion and racially polarized voting.

\textbf{Burden allocation}: Our threshold helps allocate burdens of proof. In states above 42\% minority, geographic feasibility is demonstrated empirically, so litigation focuses on prongs 2-3. In states below 37\%, proportional MM representation faces systematic geographic constraints, so plaintiffs bear heavier burden on prong 1.

\textbf{Partial answer}: We answer "Can optimization create proportional MM districts?" not "Must states create such districts?" or "Do electoral conditions warrant such districts?" Geographic feasibility is one input to legal analysis, not the complete answer.

\textbf{Complementary evidence}: Our algorithmic analysis complements traditional VRA litigation evidence (electoral analysis, expert testimony on communities of interest, historical discrimination analysis) but does not replace it.

\subsubsection{Why Algorithmic Analysis Is Valuable Despite Scope Limits}

Even addressing only prong 1, our analysis provides substantial value:

\textbf{Eliminates geographic uncertainty}: Before litigating expensive prongs 2-3 analyses, parties can assess whether geographic feasibility exists. If proportional MM representation is geographically infeasible, prongs 2-3 become moot.

\textbf{Quantifies "sufficiently large"}: Courts have struggled to define this term. Our 42\% threshold operationalizes it based on comprehensive empirical evidence.

\textbf{Reduces expert dispute}: Traditional prong 1 litigation involves expert witnesses drawing alternative maps. Our systematic optimization across 645 configurations provides objective baseline for feasibility assessment.

\textbf{Enables ex ante planning}: Legislatures and redistricting commissions can assess geographic feasibility before drawing maps, potentially avoiding litigation entirely.

\textbf{Establishes feasibility precedents}: Our 43-state analysis creates empirical foundation for evaluating prong 1 claims across diverse contexts, reducing litigation costs and increasing predictability.

\subsection{Prior Quantitative Analysis}

Limited prior work has attempted to quantify VRA feasibility thresholds:

\textbf{Legal scholarship} provides qualitative guidance but lacks systematic empirical analysis. Courts evaluate ``sufficiency'' case-by-case using expert testimony about specific geographic configurations.

\textbf{Political science} literature examines MM district creation but typically assumes feasibility rather than testing threshold conditions systematically.

\textbf{Algorithmic redistricting} research has demonstrated various optimization methods but has not explicitly studied minimum population requirements for VRA compliance.

\subsection{Gap in Literature}

\textbf{No existing work systematically identifies the minimum state-level minority percentage required for VRA compliance through algorithmic methods.} This gap leaves courts, legislatures, and advocates without quantitative tools to assess feasibility ex ante.

Our contribution fills this gap through:
\begin{itemize}
    \item Systematic testing across multiple states with varying demographics
    \item Rigorous algorithmic optimization (edge-weighted partitioning)
    \item Statistical validation of threshold patterns
    \item Geographic clustering analysis to understand spatial constraints
\end{itemize}

\subsection{Scope and Modeling Assumptions}

This study addresses a specific, bounded question: \textbf{At what state-level minority percentage can edge-weighted algorithmic redistricting create majority-minority districts in proportion to the minority population?} We must clarify what this analysis does and does not establish regarding VRA compliance.

\subsubsection{Research Question vs Legal Standard}

Our empirical analysis examines \emph{feasibility patterns}---whether optimization algorithms can achieve specified MM district targets given demographic and geographic constraints. This differs fundamentally from establishing \emph{legal requirements}:

\textbf{Research Question}: Can a state with $X$\% minority population create $Y$ MM districts (where $Y \approx X\% \times \text{total districts}$) while maintaining contiguity and reasonable compactness?

\textbf{Legal Question}: Must a state with $X$\% minority population create $Y$ MM districts to comply with VRA Section 2?

The answer to the research question informs but does not determine the answer to the legal question. Courts consider many factors beyond algorithmic feasibility when evaluating Section 2 claims.

\subsubsection{Proportional Representation as Modeling Assumption}

We adopt \textbf{proportional representation as a modeling assumption}: if a state has $X$\% minority population, we test whether optimization can create $\approx X$\% of districts as majority-minority. This assumption serves three purposes:

\begin{enumerate}
    \item \textbf{Baseline for comparison}: Proportionality provides a principled, objective target for cross-state analysis
    \item \textbf{Upper bound test}: If proportional representation is infeasible, sub-proportional representation may also face geographic constraints
    \item \textbf{Empirical regularity detection}: Systematic patterns emerge most clearly when testing consistent targets across diverse states
\end{enumerate}

\textbf{Critical clarification}: Proportional MM representation is \emph{not} a VRA legal requirement. Section 2 prohibits vote dilution where geographically compact minority communities exist; it does not mandate proportional allocation of MM districts~\citep{johnson1996}. Courts have approved plans with both more and fewer MM districts than proportional targets, depending on case-specific factors including:

\begin{itemize}
    \item Geographic distribution and clustering of minority population
    \item Communities of interest and traditional districting principles
    \item Competing constitutional requirements (equal population, contiguity)
    \item Political cohesion and racially polarized voting patterns (Gingles prongs 2-3)
\end{itemize}

Our proportionality assumption likely provides an \emph{upper bound} on VRA obligations. If states cannot achieve proportional MM representation through optimization, courts are unlikely to mandate it. Conversely, if proportional representation is readily achievable, it demonstrates geographic feasibility for at least that level of minority representation.

\subsubsection{State-Level Threshold vs District-Level Feasibility}

Our analysis identifies \emph{state-level empirical patterns}: states with $\geq$42\% minority achieve near-proportional MM representation through optimization, while states below 37\% fail to achieve proportional targets. This state-level threshold provides \textbf{contextual guidance} for evaluating VRA claims, not a bright-line legal test.

\textbf{Distinction}:
\begin{itemize}
    \item \textbf{State-level threshold (this paper)}: Empirical regularity observed across 43 states---minority percentage correlates strongly with optimization success in creating proportional MM districts
    \item \textbf{District-level feasibility (legal requirement)}: Courts evaluate whether specific districts can be drawn for specific minority communities in specific geographic configurations
\end{itemize}

VRA Section 2 operates at the district level: plaintiffs must demonstrate that a \emph{particular} minority community is sufficiently large and compact to constitute a district. Our state-level findings inform this analysis by showing that:

\begin{enumerate}
    \item \textbf{Above 42\% state minority}: Geographic feasibility for multiple MM districts is empirically demonstrated; burden shifts to defendants to explain why specific districts cannot be drawn
    \item \textbf{Below 37\% state minority}: Proportional MM representation faces systematic geographic constraints; plaintiffs face higher burden to demonstrate specific district feasibility
    \item \textbf{Borderline (37-42\%)}: Case-specific geographic analysis essential; state-level percentage alone insufficient to determine feasibility
\end{enumerate}

\subsubsection{Algorithmic Feasibility vs Legal Compliance}

Edge-weighted optimization demonstrates \emph{one approach} to MM district creation. Legal compliance involves additional considerations beyond algorithmic optimization:

\textbf{Traditional redistricting principles}: Courts balance VRA compliance against other legitimate state interests including preserving political subdivisions, maintaining communities of interest, and respecting traditional boundaries.

\textbf{Gingles prongs 2-3}: Our analysis addresses geographic compactness (prong 1) but not political cohesion (prong 2) or racially polarized voting (prong 3), which require case-specific electoral and behavioral analysis.

\textbf{Partisan considerations}: While our algorithm ignores partisan composition, real-world redistricting must consider how MM districts affect overall partisan balance and incumbency protection.

\textbf{Alternative configurations}: Our optimization finds \emph{one feasible plan}; many other configurations may exist with different tradeoffs between MM district creation and other objectives.

\subsubsection{Contribution and Scope}

This paper's contribution is \textbf{empirical and descriptive}, not normative or prescriptive:

\textbf{What we establish}: At 42\% state-level minority, edge-weighted optimization reliably creates near-proportional MM representation across diverse geographic contexts. This threshold reflects fundamental demographic-geographic constraints, not algorithmic artifacts.

\textbf{What we do not claim}: That states must create proportional MM districts, that 42\% is a legal bright-line test, or that algorithmic feasibility alone determines VRA compliance.

\textbf{Practical value}: Courts, legislatures, and redistricting commissions can use the 42\% threshold as \emph{one factor} in feasibility assessment, understanding that specific geographic configurations, community structures, and legal requirements drive ultimate compliance determinations.

Our findings transform VRA feasibility assessment from purely qualitative expert testimony to quantitatively-informed analysis, while recognizing that judicial discretion and case-specific factors remain essential to Section 2 adjudication.
